\documentclass[../main.tex]{subfiles}
\begin{document}

\chapter{ILP}
\lessoninfo{
  video={Lesson 6，動画 1--20},
  textbook={H\&P \cite{hennessy2017} §3.1，§3.4；\cite{keller1975,wall1991}},
  keywords={RAW 依存，WAW 依存，偽の依存，アーキテクチャレジスタ，物理レジスタ，レジスタリネーミング，レジスタエイリアステーブル，命令レベル並列性，理想プロセッサ，サイクル当たり命令数，プロセッサ幅，インオーダ実行，アウトオブオーダ実行},
}

分岐予測（第4章）とプレディケーション（第5章）は，制御依存によって
パイプラインが失うサイクルの大半を取り除き，CPI を理想値の~1 に近づける．
1 を下回るには複数の命令を同じサイクルで実行しなければならず，それを
妨げるのが命令間のデータ依存である．この章では，本質的な依存は RAW 依存
だけであること，それ以外の依存はレジスタリネーミングで除去されること，
そして残った依存がプログラムの命令レベル並列性---実際のプロセッサの
サイクル当たり命令数を測る基準となる上限---を定めることを示す．

\section{すべての命令を同一サイクルで}
\label{sec:l06-same-cycle}

ストールもフラッシュもしないパイプラインは毎サイクル 1 命令を完了する
ので，その CPI は~1 である（\cref{eq:pipeline-cpi}）．\source{Lesson 6，
動画 1--3；H\&P §3.1．}CPI がそれより低いということは，1 サイクルに
複数の命令が完了するということである．極端な場合，プログラムの $N$ 個の
命令すべてが同じサイクルでフェッチされ，次のサイクルで一斉にデコード
され，その次のサイクルで一斉に実行され，以下同様に進む
（\cref{fig:l06-same-cycle}）．5 段プロセッサはプログラム全体を
5 サイクルで完了し，CPI は $5/N$ となる．これは長いプログラムでは~0 に
近づく．

実際のプログラムはこのようには実行できない．命令が他の命令の結果を使う
からである．この章の前半を通じて用いる例として，次の 5 命令を考える．%
\begin{marginfigure}
  \resizebox{\linewidth}{!}{% Thought experiment: all N instructions pass through the five stages in
% the same five cycles.
\begin{tikzpicture}[x=0.9cm,y=0.55cm, font=\footnotesize\sffamily,
  st/.style={draw, minimum width=0.82cm, minimum height=0.46cm, inner sep=0pt, fill=ln-box}]
  \foreach \c in {1,...,5} { \node at (\c,0.8) {\c}; }
  \node[anchor=east] at (0.5,0.8) {\iflangja{サイクル}{cycle}};
  \foreach \r/\name in {0/I1,1/I2,2/I3,4/I$N$} {
    \node[anchor=east] at (0.5,-\r) {\name};
    \foreach \s/\lab in {1/IF,2/ID,3/EX,4/MEM,5/WB} {
      \node[st] at (\s,-\r) {\lab};
    }
  }
  \foreach \s in {1,...,5} { \node at (\s,-3) {$\vdots$}; }
\end{tikzpicture}
}
  \caption{理想的な場合．$N$ 個の命令すべてが同じ 5 サイクルで各ステージ
    を通過する．}
  \label{fig:l06-same-cycle}
\end{marginfigure}
\begin{lstlisting}[language={[x86masm]Assembler}, numbers=none]
I1:  ADD R3, R1, R2     ; R3 <- R1 + R2
I2:  SUB R5, R3, R4     ; R5 <- R3 - R4
I3:  AND R6, R1, R4     ; R6 <- R1 AND R4
I4:  OR  R4, R2, R7     ; R4 <- R2 OR R7
I5:  XOR R5, R1, R7     ; R5 <- R1 XOR R7
\end{lstlisting}
5 命令すべてが同じサイクルで実行ステージにあるとき，I2 は，I1 がまだ R3
を計算している間に $\text{R3} - \text{R4}$ を計算するので，R3 の古い値を
使ってしまう．フォワーディング（\cref{fig:forwarding}）ではこれを正せ
ない．フォワーディングは，あるサイクルで計算された結果をより後のサイクル
で実行される命令に届けるものであるが，ここでは結果とその使用が同じ
サイクルに属するからである．したがって，プロセッサがどこまで理想に
近づけるかは，プログラムのデータ依存によって決まる．

\section{同一サイクル実行における依存}
\label{sec:l06-dependences}

I1 に対する I2 の RAW 依存は守らなければならない．I2 を実行できるのは
早くても I1 の 1 サイクル後，すなわち I1 の結果をフォワーディングで
受け取れるときである．\source{Lesson 6，動画 4--6；H\&P §3.1．}I3，I4，
I5 は，先行する命令が書き込むレジスタを読まないので，依然として I1 と
同じサイクルで実行できる．\cref{fig:l06-dependences} はその結果の
タイミングを示す．I2 は 1 サイクル待ち，I3 から I5 は，プログラム上では
I2 の後にあるにもかかわらず I2 より先に実行される．したがって RAW 依存は，
幅に制限がなくても CPI をどこまで下げられるかを制約する．すべての命令が
直前の命令の結果を読むなら，毎サイクル 1 命令しか完了せず，CPI は~1 の
ままである．

\begin{figure}[htbp]
  \centering
  % Same-cycle execution of the five-instruction program: the RAW dependence
% delays I2, which then overwrites the later write of I5 to R5 (WAW).
\begin{tikzpicture}[x=1.2cm,y=0.9cm, font=\footnotesize\sffamily, >={Stealth[length=1.8mm]},
  st/.style={draw, minimum width=0.95cm, minimum height=0.5cm, inner sep=0pt, fill=ln-box},
  bad/.style={st, fill=ln-caution!15, draw=ln-caution},
  stall/.style={draw=ln-rule, dashed, minimum width=0.95cm, minimum height=0.5cm,
    inner sep=0pt, font=\scriptsize\sffamily, text=ln-muted}]
  \foreach \c in {1,...,6} { \node at (\c,0.7) {\c}; }
  \node[anchor=east] at (0.45,0.7) {\iflangja{サイクル}{cycle}};
  \foreach \r/\txt in {1/{ADD R3,R1,R2},2/{SUB R5,R3,R4},3/{AND R6,R1,R4},
                       4/{OR~~R4,R2,R7},5/{XOR R5,R1,R7}} {
    \node[anchor=east] at (0.45,1-\r) {I\r\enspace\texttt{\txt}};
  }
  % I1, I3, I4: undelayed
  \foreach \s/\lab in {1/IF,2/ID,4/MEM,5/WB} { \node[st] at (\s,0) {\lab}; }
  \node[st] (e1) at (3,0) {EX};
  \foreach \r in {3,4} {
    \foreach \s/\lab in {1/IF,2/ID,3/EX,4/MEM,5/WB} { \node[st] at (\s,1-\r) {\lab}; }
  }
  % I2: waits one cycle for R3
  \node[st] at (1,-1) {IF};
  \node[st] at (2,-1) {ID};
  \node[stall] at (3,-1) {\iflangja{待ち}{wait}};
  \node[st] (e2) at (4,-1) {EX};
  \node[st] at (5,-1) {MEM};
  \node[bad] (w2) at (6,-1) {WB};
  % I5: independent, but writes R5 before I2 does
  \foreach \s/\lab in {1/IF,2/ID,3/EX,4/MEM} { \node[st] at (\s,-4) {\lab}; }
  \node[bad] (w5) at (5,-4) {WB};
  % forwarding of R3 from I1 to I2
  \draw[->, thick, ln-accent] (e1.south) -- (3,-0.5)
    -- node[fill=white, inner sep=0.6pt, font=\scriptsize\sffamily] {R3} (4,-0.5)
    -- (e2.north);
  % WAW violation
  \draw[->, dashed, thick, ln-caution] (w2.south) -- (w5.north east);
  \node[anchor=west, font=\scriptsize\sffamily, text=ln-caution, align=left] at (6.0,-2.5)
    {\iflangja{R5 を\\上書き}{overwrites\\R5}};
\end{tikzpicture}

  \caption{例の 5 命令の同一サイクル実行．R3 に関する RAW 依存が I2 を
    1 サイクル遅らせる．その結果 I2 は I5 より後に R5 に書き込み，両者の
    WAW 依存に違反する．}
  \label{fig:l06-dependences}
\end{figure}

先行する命令を追い越して命令を実行すると，I2 と I5 の間の WAW 依存が
破られる．%
\tip{名前の後半は\emph{先行する}命令の動作を表す．write-after-read なら
先行命令が読む側である．}両者はともに R5 に書き込み，プログラム順では I5 が最後に書く
ので，これらの命令の後で R5 は I5 の結果を保持していなければならない．
R5 を読む後続の命令はすべてその値を前提としている．しかし
\cref{fig:l06-dependences} では，I5 がサイクル~5 で R5 に書き込み，
遅延した I2 がサイクル~6 でそれを上書きする．正しい値を保つには，I5 の
計算は I2 から何も必要としないにもかかわらず，I5 も I2 の書き込みの後
まで遅らせなければならない．

WAR 依存も同じように危うくなる．I4 は I2 が読む R4 に書き込み，プログラム
順では I2 は R4 の古い値を使わなければならない．I2 が遅延すると，I2 の
実行時には I4 が生成した新しい値がすでに存在しており，I2 がそれを使わない
ようにしなければならない．このように，1つの RAW 遅延は波及する．遅延した
命令が読み書きするレジスタに書き込む後続の命令は，すべて同様に遅らせ
なければならない．そうしてはじめて，各命令はプログラム中の自分の位置に
対応するレジスタ値を見ることができる．

\section{偽の依存の除去}
\label{sec:l06-false}

第3章の 3 種類のデータ依存は性質が異なる．\source{Lesson 6，動画 7--8；
H\&P §3.1．}RAW 依存では，後の命令が先の命令の計算した値を使う．データは
実際に一方の命令から他方へ流れるのであり，どのようなハードウェアでも，
読む側の命令を書く側の命令より先に実行させることはできない．%
\index{RAW いぞん@RAW 依存}WAW 依存と WAR 依存%
\index{WAW いぞん@WAW 依存}\index{WAR いぞん@WAR 依存}は偽の
依存\index{ぎのいぞん@偽の依存}である．2つの命令の間で値は受け渡され
ない．%
\caution{H\&P はこれらを\emph{名前依存}と呼ぶ．WAR は\emph{逆依存}，
WAW は\emph{出力依存}である．}これらの依存は，プログラムが異なる値を同じ
レジスタに格納するという
だけの理由で存在する．それは，レジスタの数に限りがあること，そして繰り
返し実行される命令---例えばループの各反復で実行される命令---が常に同じ
レジスタに書き込むことの帰結である．この例では，I2 と I5 の結果は，
たまたま R5 という名前が付いた無関係な値である．

したがって，異なる値を異なる場所に保持すれば偽の依存は消える．その方法の
1つは，レジスタの複数のバージョンを保持することである．I2 と I5 がともに
R5 に書き込むとき，プロセッサは一方で他方を上書きするのではなく両方の
結果を格納し，R5 を読む各命令はプログラム順によって選ばれるバージョンを
使う．I2 と I5 の間にある命令は I2 の結果を，I5 の後にある命令は I5 の
結果を使う．2つの書き込みが実際に起こる順序はもはや問題にならず，I5 は
I2 が存在しない場合と同じだけ早く実行できる．同様に，I2 は，I4 が新しい
バージョンをいつ生成するかにかかわらず，R4 の古いバージョンを読む．

難しいのはその管理である．読み出しのたびに，プロセッサは正しい
バージョン，すなわちプログラム順で最も近い先行する書き込み命令が生成した
バージョンを見つけなければならず，またバージョンが不要になったことを
認識しなければならない．レジスタリネーミングはこの管理を系統的に行う．

\section{レジスタリネーミング}
\label{sec:l06-renaming}

リネーミングは，プログラムが名前で指定するレジスタと，プロセッサがその
値を保持する記憶領域とを分離する．\source{Lesson 6，動画 9--11；H\&P
§3.1，§3.4；\cite{keller1975}．}

\begin{definition}[アーキテクチャレジスタと物理レジスタ]
  \label{def:l06-registers}
  命令セットで定義され，したがってプログラム中で名前で指定されるレジスタ
  （例えば R0 から R31）を
  \term[あーきてくちゃれじすた]{アーキテクチャレジスタ}[architectural register]%
  という．プロセッサがレジスタの値を
  実際に保持する記憶位置を\term[ぶつりれじすた]{物理レジスタ}[physical
  register]という．リネーミングを行うプロセッサは，アーキテクチャ
  レジスタよりはるかに多くの物理レジスタを持つ．
\end{definition}

\begin{definition}[レジスタリネーミング]
  \label{def:l06-renaming}
  各命令をデコード時に書き換え，アーキテクチャレジスタの代わりに物理
  レジスタを読み書きさせ，結果ごとに新たに割り当てた物理レジスタを用いる
  ようにすることを\term[れじすたりねーみんぐ]{レジスタリネーミング}[register
  renaming]という．
\end{definition}

\begin{definition}[レジスタエイリアステーブル]
  \label{def:l06-rat}
  \term[れじすたえいりあすてーぶる]{レジスタエイリアステーブル}[register alias table]%
  \index{RAT|see{レジスタエイリアステーブル}}（RAT）は，アーキテクチャ
  レジスタごとに 1つのエントリを持つ．各エントリは，それまでにリネーミング
  された命令の中でプログラム順に最も新しい，そのアーキテクチャレジスタの
  値を保持する物理レジスタを示す．
\end{definition}

各命令は 3つのステップでリネーミングされる（\cref{fig:l06-rat}）．%
\margin{講義では RAT を \emph{register allocation table} の略としている．
どちらの名前も同じテーブルを指す．}
\begin{enumerate}
  \item 各ソースレジスタについて，その RAT エントリを読む．リネーミング後
    の命令はその物理レジスタを読む．
  \item 結果のために，空いている物理レジスタ，すなわちまだ必要な値を保持
    していない物理レジスタを割り当てる．リネーミング後の命令はそこに書き
    込む．
  \item 宛先レジスタの RAT エントリを新しい物理レジスタで上書きする．
    これにより，このアーキテクチャレジスタを読む後続の命令は新しい結果を
    使う．
\end{enumerate}
ソースの参照は，宛先のエントリを上書きする前に行わなければならない．
そうしないと，\texttt{ADD R1, R1, R2} のような命令が R1 の以前の値ではなく
自分自身の結果を読んでしまう．\margin{実際の設計でリネーミングがアウト
オブオーダ実行とどのように組み合わされるかは，第7章と第8章で示す．}

\begin{figure}[htbp]
  \centering
  % Renaming one instruction: sources are looked up in the RAT, the result
% gets a free physical register, and the RAT entry of the destination is
% updated to it.
\begin{tikzpicture}[x=1cm,y=1cm, font=\footnotesize\sffamily, >={Stealth[length=1.8mm]},
  fld/.style={draw, fill=ln-box, minimum width=1.1cm, minimum height=0.5cm, inner sep=0pt},
  arch/.style={draw, fill=ln-box, minimum width=0.8cm, minimum height=0.5cm, inner sep=0pt},
  phys/.style={draw, fill=white, minimum width=0.8cm, minimum height=0.5cm, inner sep=0pt}]
  % decoded instruction
  \node[fld] (s1) at (0,0)    {R1};
  \node[fld] (s2) at (0,-0.5) {R2};
  \node[fld] (d)  at (0,-1.0) {R3};
  \foreach \n/\lab in {s1/src$_1$,s2/src$_2$,d/dst} {
    \node[anchor=east, font=\scriptsize, text=ln-muted] at (\n.west) {\lab};
  }
  \node[anchor=south] at (0,0.3) {\texttt{ADD R3, R1, R2}};
  % RAT
  \foreach \i in {1,...,7} {
    \node[arch] (a\i) at (3.0,{-(\i-1)*0.5}) {R\i};
  }
  \foreach \i in {1,2,4,5,6,7} {
    \node[phys] (p\i) at (3.8,{-(\i-1)*0.5}) {P\i};
  }
  \node[phys, fill=ln-accent!15] (p3) at (3.8,-1.0) {P8};
  \node[anchor=north, inner sep=1pt] at (3.4,-3.27) {$\vdots$};
  \node[anchor=north, font=\footnotesize\bfseries\sffamily] at (3.4,-3.75) {RAT};
  % renamed instruction
  \node[fld] (r1) at (6.8,0)    {P1};
  \node[fld] (r2) at (6.8,-0.5) {P2};
  \node[fld] (rd) at (6.8,-1.0) {P8};
  \node[anchor=south] at (6.8,0.3) {\texttt{ADD P8, P1, P2}};
  % lookups
  \draw[->] (s1.east) -- (a1.west);
  \draw[->] (s2.east) -- (a2.west);
  \draw[->] (p1.east) -- (r1.west);
  \draw[->] (p2.east) -- (r2.west);
  % destination
  \draw[->, dashed] (d.east) -- (a3.west)
    node[midway, below, font=\scriptsize] {\iflangja{エントリ選択}{select entry}};
  \node[draw, fill=ln-box, minimum height=0.5cm, inner sep=3pt] (free) at (6.8,-2.4)
    {\iflangja{空き}{free}: P8, P9, P10, \ldots};
  \draw[->, thick, ln-accent] (free.north) -- (rd.south);
  \draw[->, thick, ln-accent] (rd.west) -- (p3.east)
    node[midway, below, font=\scriptsize, text=ln-accent] {\iflangja{更新}{update}};
\end{tikzpicture}

  \caption{\texttt{ADD R3, R1, R2} のリネーミング．R1 と R2 を RAT で参照
    し，結果には空いている物理レジスタ P8 を与え，P3 を保持していた R3 の
    RAT エントリを P8 で上書きする．}
  \label{fig:l06-rat}
\end{figure}

物理レジスタは，\cref{sec:l06-false}で述べたバージョンにほかならない．
書き込みのたびに新しいバージョンが作られ，RAT はプログラム順でどの
バージョンが現在のものかを記録するので，読み出す命令は自動的に，最も
近い先行する書き込み命令が生成したバージョンを使う．%
\sidenote{リネーミングは命令そのものを肩代わりすることもある．Ivy Bridge
以降の Intel のコアは，レジスタ間のムーブの大半を RAT だけで処理し，宛先の
エントリをソースの物理レジスタに向ける．このムーブは実行ユニットを使わず，
レイテンシもない（Agner Fog，マイクロアーキテクチャマニュアル）．}

\begin{example}
  \label{ex:l06-rename}
  例の 5 命令をリネーミングする．I1 の前には，RAT は例で使うレジスタ R1 から
  R7 を P1 から P7 に，その他のアーキテクチャレジスタを例で使わない物理
  レジスタに対応付けている．P8，P9，\ldots\ は空いており，番号の小さい順に
  割り当てられる．%
  \margin{実際の規模：Skylake は物理レジスタを整数用に 180 個，ベクタ用に
  168 個持つ．アーキテクチャレジスタはそれぞれ 16 個である（Agner Fog，
  マイクロアーキテクチャマニュアル）．}\cref{tab:l06-rename} は，リネーミング後の各命令と，
  その命令の後の RAT エントリを示す．I2 が I1 の結果である P8 を読むのは，
  I1 がすでに R3 のエントリを更新しているからである．I3 は，I4 が R4 を
  P11 に対応付ける前に，古い R4 である P4 を読む．R5 への 2つの書き込みは
  P9 と P12 に向かい，最後に RAT は R5 を P12 に対応付けるので，後続の命令
  はプログラム順が要求するとおり I5 の結果を読む．
\end{example}

\begin{table}[htbp]
  \centering\small
  \caption{例の 5 命令のリネーミング．命令によって変更されたエントリを太字で
    示す．R1，R2，R7 は終始 P1，P2，P7 のままである．}
  \label{tab:l06-rename}
  \begin{tabular}{@{}lllcccc@{}}
    \toprule
    & 命令 & リネーミング後 & R3 & R4 & R5 & R6 \\
    \midrule
    & （初期状態） & & P3 & P4 & P5 & P6 \\
    I1 & \texttt{ADD R3, R1, R2} & \texttt{ADD P8, P1, P2}
       & \textbf{P8} & P4 & P5 & P6 \\
    I2 & \texttt{SUB R5, R3, R4} & \texttt{SUB P9, P8, P4}
       & P8 & P4 & \textbf{P9} & P6 \\
    I3 & \texttt{AND R6, R1, R4} & \texttt{AND P10, P1, P4}
       & P8 & P4 & P9 & \textbf{P10} \\
    I4 & \texttt{OR~~R4, R2, R7} & \texttt{OR~~P11, P2, P7}
       & P8 & \textbf{P11} & P9 & P10 \\
    I5 & \texttt{XOR R5, R1, R7} & \texttt{XOR P12, P1, P7}
       & P8 & P11 & \textbf{P12} & P10 \\
    \bottomrule
  \end{tabular}
\end{table}

\section{リネーミング後の依存}
\label{sec:l06-after-renaming}

リネーミングされたプログラムでは，すべての結果がそれぞれ専用の物理
レジスタに書き込まれる．\source{Lesson 6，動画 12；H\&P §3.4．}2つの命令が
同じ位置に書き込むことはなく，先行する命令が読む位置に書き込む命令もない
ので，レジスタアクセス間の WAW 依存と WAR 依存はもはや存在しない．%
\margin{メモリ位置を介した依存はレジスタリネーミングの影響を受けない．
それは第9章で扱う．}一方，すべての読み出しは，元のプログラムでその値を
生成した命令の結果を依然として参照するので，RAW 依存はすべて保たれる．
\cref{tab:l06-rename} で残る依存は，P8 を介した I1 に対する I2 の依存
だけである．

同一サイクル実行が守らなければならないのは，今やこの 1つの依存だけで
ある．I1，I3，I4，I5 は最初のサイクルで，I2 は 2 番目のサイクルで実行
され，I1 の結果は I2 にフォワーディングされる．I2 と I5 はそれぞれ P9 と
P12 に書き込むので書き込みの順序は無関係であり，I2 が読む P4 に I4 は
触れない．5 命令は 2 サイクルで実行され，後続の読み出し命令はすべて RAT
を通じて正しい値を見つける．

\begin{keypoint}
  RAW 依存は命令間の唯一の真の依存である．レジスタリネーミングはレジスタ
  に関するすべての WAW 依存と WAR 依存を除去するので，命令の並列実行に
  残る制約はデータの流れそのものだけである．
\end{keypoint}

\section{命令レベル並列性}
\label{sec:l06-ilp}

偽の依存を除去すれば，%
\source{Lesson 6，動画 13--15；H\&P 第3章；\cite{wall1991}．}%
プログラムの実行に必要なサイクル数はその RAW 依存だけで決まる．この数は
プログラム自体の性質であり，特定のプロセッサの性質ではない．これは仮想的
な\term[りそうぷろせっさ]{理想プロセッサ}[ideal processor]上で測定される．
理想プロセッサは次の性質を持つ．
\begin{itemize}
  \item レジスタをリネーミングするので，偽の依存は存在しない．
  \item 各命令を，それが読む値がすべて得られる最初のサイクルで実行する．
    それは，先行する命令に依存しなければ最初のサイクル，依存すれば依存先
    の命令のうち最も遅いものの次のサイクルである．
  \item 1 サイクルに実行する命令数に制限がない．
  \item ハードウェアに制限がないので，使用中の実行ユニットを待つ命令は
    ない．すなわち構造依存は存在しない．
  \item 各命令を，演算の種類によらずちょうど 1 サイクルで実行する．
\end{itemize}

\begin{definition}[命令レベル並列性]
  \label{def:l06-ilp}
  $N$ 命令を実行するプログラムの
  \term[めいれいれべるへいれつせい]{命令レベル並列性}[instruction-level parallelism]%
  \index{ILP|see{命令レベル並列性}}（ILP）を次式で定義する．
  \begin{equation}
    \text{ILP} \;=\; \frac{N}{C} ,
    \label{eq:l06-ilp}
  \end{equation}
  ここで $C$ は，理想プロセッサがそれらの命令を実行するサイクル数である．
\end{definition}

理想プロセッサが命令 $I$ を実行するサイクル $c(I)$ は，RAW 依存から直接
求まる．
\begin{equation}
  c(I) \;=\; 1 + \max_{J \in D(I)} c(J) ,
  \qquad C = \max_{I} c(I) ,
  \label{eq:l06-cycle}
\end{equation}
ここで $D(I)$ は $I$ が結果を読む命令の集合であり，空集合に対する最大値は~0
とする．したがって $C$ は，プログラム中の RAW 依存の最長の連鎖の長さ
である．どのプロセッサも，命令が必要とする値が存在する前にその命令を実行
することはできないので，ILP は，任意のプロセッサがこのプログラムを実行
できる速度（サイクル当たりの命令数）の上限である（\cref{sec:l06-ipc}）．%
\tip{紙の上でリネーミングを行う必要はない．各レジスタの読み出しについて，
最も近い\emph{先行する}書き込み命令を探せばよい．後の書き込みは関係
しない．}

\needspace{12\baselineskip}
\begin{example}
  \label{ex:l06-ilp}
  次の 8 命令の ILP を求める．
  \begin{lstlisting}[language={[x86masm]Assembler}, numbers=none]
I1:  ADD R1, R2, R3
I2:  MUL R4, R1, R5    ; R1 は I1 から，R5 は後で書かれる
I3:  SUB R6, R2, R7    ; R2，R7 は後でしか書かれない
I4:  ADD R1, R6, R8    ; R6 は I3 から
I5:  AND R9, R1, R8    ; R1 は I1 からではなく I4 から
I6:  OR  R2, R8, R3
I7:  XOR R7, R4, R2    ; R4 は I2 から，R2 は I6 から
I8:  ADD R5, R3, R8
\end{lstlisting}
  I1，I3，I6，I8 は，それより前に書き込まれていないレジスタだけを読む
  ので，サイクル~1 で実行される．I2 と I4 はサイクル~1 の命令に依存し，
  サイクル~2 で実行される．I5 と I7 はそれぞれサイクル~2 の命令に依存し，
  サイクル~3 で実行される（\cref{fig:l06-ilp-graph}）．よって $C=3$，
  $\text{ILP} = 8/3 \approx 2.67$ である．%
  \sidenote{実際のプログラムでもこれより大きくはない．Wall は，現実的な
  ものから実現不可能なほど理想的なもの---完全な分岐予測，無制限のリネー
  ミング，完全なエイリアス解析---までのモデルでシミュレーションを行い，
  平均の並列度が~7 を超えることはまれで，5 程度が多いことを示した
  \cite{wall1991}．}ここには 2つの落とし穴が見られる．
  R1 には I1 と I4 の両方が書き込み，I5 が読むのは I4 の値である．I1 を
  生成元とみなすと，I5 を誤ってサイクル~2 に置いてしまう．また I3 は R2 と
  R7 を読むが，I6 と I7 がそれらに書き込むのは後になってからである．これら
  は WAR 依存であり，リネーミングで除去されるので，I3 を遅らせない．
\end{example}

\begin{figure}[htbp]
  \centering
  % RAW dependences of the ILP example, instructions placed in the cycle in
% which the ideal processor executes them.
\begin{tikzpicture}[x=1cm,y=1cm, font=\footnotesize\sffamily, >={Stealth[length=1.8mm]},
  ins/.style={draw, fill=ln-box, rounded corners=2pt, align=center,
    minimum width=2.4cm, minimum height=0.75cm, inner sep=1pt,
    font=\scriptsize\sffamily},
  dep/.style={->, thick, ln-accent},
  reg/.style={font=\scriptsize\sffamily, fill=white, inner sep=1pt}]
  \foreach \c/\x in {1/0,2/3.9,3/7.8} {
    \node[font=\footnotesize\bfseries\sffamily, text=ln-muted] at (\x,0.75)
      {\iflangja{サイクル \c}{cycle \c}};
  }
  \draw[ln-rule] (1.95,0.95) -- (1.95,-3.4);
  \draw[ln-rule] (5.85,0.95) -- (5.85,-3.4);
  % cycle 1
  \node[ins] (i1) at (0,0)    {I1\\\texttt{ADD R1,R2,R3}};
  \node[ins] (i6) at (0,-1)   {I6\\\texttt{OR~~R2,R8,R3}};
  \node[ins] (i3) at (0,-2)   {I3\\\texttt{SUB R6,R2,R7}};
  \node[ins] (i8) at (0,-3)   {I8\\\texttt{ADD R5,R3,R8}};
  % cycle 2
  \node[ins] (i2) at (3.9,0)  {I2\\\texttt{MUL R4,R1,R5}};
  \node[ins] (i4) at (3.9,-2) {I4\\\texttt{ADD R1,R6,R8}};
  % cycle 3
  \node[ins] (i7) at (7.8,-0.5) {I7\\\texttt{XOR R7,R4,R2}};
  \node[ins] (i5) at (7.8,-2)   {I5\\\texttt{AND R9,R1,R8}};
  % RAW dependences, labelled with the register
  \draw[dep] (i1) -- node[reg] {R1} (i2);
  \draw[dep] (i3) -- node[reg] {R6} (i4);
  \draw[dep] (i4) -- node[reg] {R1} (i5);
  \draw[dep] (i2.east) -- node[reg] {R4} (i7.north west);
  \draw[dep] (i6.east) -- node[reg, pos=0.3] {R2} (i7.south west);
\end{tikzpicture}

  \caption{\cref{ex:l06-ilp} の RAW 依存．矢印には値を運ぶレジスタを
    付した．各命令は，理想プロセッサがそれを実行するサイクルに置かれて
    いる．最長の連鎖の長さは~3 である．}
  \label{fig:l06-ilp-graph}
\end{figure}

\section{構造依存と制御依存}
\label{sec:l06-control}

構造依存，すなわちある命令が別の命令の使用しているハードウェアを待つ
依存は，準備のできたすべての命令に十分なハードウェアを持つ理想プロセッサ
では生じない．\source{Lesson 6，動画 16；H\&P §3.1．}%
\caution{H\&P はこれを構造\emph{ハザード}と呼ぶ．プログラムではなく
ハードウェアに由来するので，本来は依存ではない．}しかし，分岐は制御
依存を加える．分岐の後の命令が実行されるかどうかは，分岐の結果に依存する．
ILP を定義する理想プロセッサは，すべての分岐を完全に予測するとも仮定
する．したがって，どの命令が実行されるかを最初から知っており，分岐の後の
命令は分岐の実行を待たない．分岐自体は他の命令と同じく 1つの命令であり，
$N$ に数えられ，そのオペランドを生成する命令の 1 サイクル後に実行される．
分岐の結果によって飛ばされる命令は実行されず，数えられない．%
\margin{プレディケーション（第5章）は制御依存そのものを取り除く．ただし
両方の経路を実行する代償を伴う．}

\begin{example}
  \label{ex:l06-control}
  次の断片では分岐が成立するので，I4 は実行されない．
  \begin{lstlisting}[language={[x86masm]Assembler}, numbers=none]
I1:        ADD R1, R2, R3     ; サイクル 1
I2:        SUB R4, R1, R5     ; サイクル 2: R1 は I1 から
I3:        BEQ R4, R6, Skip   ; サイクル 3: R4 は I2 から
I4:        ADD R7, R8, R9     ; 実行されない
I5:  Skip: XOR R8, R2, R6     ; サイクル 1
I6:        AND R9, R8, R1     ; サイクル 2: R8 は I5 から
\end{lstlisting}
  分岐を完全に予測すれば，分岐がサイクル~3 になってようやく実行される
  にもかかわらず，I5 はサイクル~1 で実行される．5 命令が 3 サイクルで
  実行されるので，$\text{ILP} = 5/3 \approx 1.67$ である．I5 と I6 が分岐
  を待たなければならないとすると，それらはサイクル 4 と~5 で実行され，
  ILP は~1 に下がる．
\end{example}

\section{ILP と IPC}
\label{sec:l06-ipc}

ILP は理想プロセッサ上でのプログラムを特徴づける量であり，実際の
プロセッサを特徴づけるのは，それが実際に達成する速度である．
\source{Lesson 6，動画 17--20；H\&P 第3章．}

\begin{definition}[サイクル当たり命令数]
  \label{def:l06-ipc}
  プロセッサが 1 サイクル当たりに完了する命令の平均数を，そのプロセッサの
  \term[さいくるあたりめいれいすう]{サイクル当たり命令数}[instructions per cycle]%
  \index{IPC|see{サイクル当たり命令数}}（IPC）という．IPC は性能の鉄則
  （\cref{eq:iron-law}）の CPI の逆数であり，$\text{IPC} =
  1/\text{CPI}$ である．%
  \caution{プログラム全体の IPC は総命令数を総サイクル数で割った値であり，
  各区間の IPC の平均ではない．}
\end{definition}

すべてのプロセッサとプログラムについて $\text{IPC} \le \text{ILP}$ が
成り立つ．理想プロセッサは，各命令を依存が許す限り早く実行するからで
ある．実際のプロセッサは，いくつかの理由で理想に届かない．
\begin{description}
  \item[幅] 1 サイクルに実行できる命令数は有限の $w$ に限られる．
    \emph{幅の狭い}プロセッサは 1つか 2つ，\emph{幅の広い}プロセッサは
    4つ以上の命令を実行する．プログラムによらず $\text{IPC} \le w$ で
    ある．%
    \margin{実際の幅：Skylake は 1 クロック 4 命令，Ice Lake は~5 命令を
    狙った設計である（Agner Fog，マイクロアーキテクチャマニュアル）．}
  \item[実行ユニット] 加算器，乗算器などの実行ユニットの数は限られており，
    乗算のような演算には複数サイクルかかることがある．必要な種類のユニット
    が使用中である命令は，入力が揃っていても待つ．これは構造依存であり，
    理想プロセッサにはないものである．
  \item[実行順序] \term[いんおーだじっこう]{インオーダ実行}[in-order
    execution]を行うプロセッサは，プログラム上で先行する命令より先に命令を
    実行することは決してない．連続する独立な命令は
    同じサイクルで実行できるが，入力を待たなければならない命令は，その
    後ろにあるすべての命令を引き止める．\emph{アウトオブオーダ}%
    \index{あうとおぶおーだじっこう@アウトオブオーダ実行}実行を行う
    プロセッサは，入力が揃った任意の命令を実行できる（第7章）．
\end{description}
アウトオブオーダプロセッサがそうするにはリネーミングが必要である．
リネーミングがなければ，\cref{sec:l06-dependences}の WAW 依存と WAR 依存
も守らなければならない．インオーダプロセッサはレジスタをプログラム順に
読み書きするので，リネーミングを必要としない．

\begin{example}
  \label{ex:l06-ipc}
  \cref{ex:l06-ilp} のプログラムを 4つのプロセッサで実行する．いずれの
  プロセッサも各命令を 1 サイクルで実行し，あらゆる種類の実行ユニットを
  十分に持つ．アウトオブオーダのプロセッサはレジスタをリネーミングし，
  準備のできた命令が幅を超えるときは古い命令から実行する．
  \cref{tab:l06-ipc} は各サイクルで実行される命令を示す．インオーダの
  プロセッサはサイクル~1 で I1 しか実行しない．I2 が R1 を待たなければ
  ならず，I3 から I8 をせき止めるからである．同じせき止めが I4（R6 待ち），
  I5（I4 からの R1 待ち），I7（R2 待ち）でも繰り返されるので，これらの
  プロセッサは幅によらず 5 サイクルを要する．幅 2 のアウトオブオーダ
  プロセッサは毎サイクル準備のできた 2 命令を実行し，4 サイクルを要する．
  ILP に到達するのは，幅 4 のアウトオブオーダプロセッサだけである．
\end{example}

\begin{table}[htbp]
  \centering\small
  \caption{\cref{ex:l06-ilp} の命令のうち，各サイクルで実行されるもの．
    理想プロセッサのスケジュールは，幅 4 のアウトオブオーダプロセッサの
    ものと等しい．}
  \label{tab:l06-ipc}
  \begin{tabular}{@{}lcccccc@{}}
    \toprule
    & \multicolumn{5}{c}{サイクル} & \\
    \cmidrule(lr){2-6}
    プロセッサ & 1 & 2 & 3 & 4 & 5 & IPC \\
    \midrule
    インオーダ，幅 2       & 1          & 2, 3 & 4    & 5, 6 & 7, 8 & 1.60 \\
    インオーダ，幅 4       & 1          & 2, 3 & 4    & 5, 6 & 7, 8 & 1.60 \\
    アウトオブオーダ，幅 2 & 1, 3       & 2, 4 & 5, 6 & 7, 8 &      & 2.00 \\
    アウトオブオーダ，幅 4 & 1, 3, 6, 8 & 2, 4 & 5, 7 &      &      & 2.67 \\
    \bottomrule
  \end{tabular}
\end{table}

この例では，2つのインオーダプロセッサはいずれも順序に制約されており，幅を
2 から 4 に広げても何も得られない．アウトオブオーダ実行は，幅 2 の
プロセッサをその幅まで（1.60 から $w = 2$ へ）しか引き上げないが，幅 4 の
プロセッサは ILP まで（1.60 から 2.67 へ）引き上げる．一般には
（\cref{tab:l06-limits}），幅の狭いプロセッサを主に制約するのは幅である
（$w = 1$ ならどちらの順序でも $\text{IPC} \le 1$）ので，アウトオブオーダ
実行による IPC の向上は限定的である．幅の広いインオーダプロセッサでは，
待っている命令が後ろの命令をせき止め，幅の大半が使われないままになる．
したがって，幅の広いプロセッサはその幅を活かすためにアウトオブオーダ実行
を必要とし，アウトオブオーダ実行は幅の広いプロセッサではじめて十分に
効果を発揮する．

\begin{table}[htbp]
  \centering\small
  \caption{プログラムの ILP に対して，プロセッサの IPC を制限するもの．}
  \label{tab:l06-limits}
  \begin{tabular}{@{}lll@{}}
    \toprule
    & インオーダ & アウトオブオーダ \\
    \midrule
    幅が狭い（幅 1--2）   & 主に幅 & 主に幅 \\
    幅が広い（幅 4 以上） & 主に順序（幅は使われない） & わずか（IPC は ILP に近づく） \\
    \bottomrule
  \end{tabular}
\end{table}

\begin{keypoint}[章のまとめ]
  多くの命令を同じサイクルで実行すれば CPI は 1 を大きく下回るはずだが，
  データ依存がそれを妨げる．真の依存は RAW 依存だけであり，WAW 依存と
  WAR 依存はレジスタ名の再利用から生じる．レジスタエイリアステーブルに
  基づくレジスタリネーミングは，アーキテクチャレジスタを新たに割り当てた
  物理レジスタに対応付け，レジスタアクセス間の偽の依存をすべて除去する．
  プログラムの ILP（\cref{eq:l06-ilp}）は，リネーミングを行い，幅に制限が
  なく，構造依存がなく，分岐を完全に予測する理想プロセッサのサイクル当たり
  命令数，すなわち命令数を RAW 依存の最長の連鎖の長さで割った値である．
  これはプログラムの性質であり，あらゆるプロセッサの IPC の上限となる．
  実際のプロセッサは，幅，実行ユニット，インオーダ実行に制約されてこれを
  下回る．ILP に近づくには幅の広いアウトオブオーダプロセッサが必要である．%
  \margin{どちらも同じスマートフォンに載る．Arm の Cortex-A55 は幅 2 の
  インオーダ，Cortex-A78 は幅 4 のアウトオブオーダである．}
  次章では，リネーミングを伴って命令をアウトオブオーダに実行する Tomasulo
  のアルゴリズムを紹介する．
\end{keypoint}

\end{document}
